\medskip

On considère la fonction $f$ définie sur $\R$ par 
\[f(x) = \ln \left(1 + \text{e}^{-x}\right),\]

où ln désigne la fonction logarithme népérien.
 
On note $\mathcal{C}$ sa courbe représentative dans un repère orthonormé \Oij. 

La courbe $\mathcal{C}$ est tracée ci-dessous.

\begin{center}
\psset{unit=2.25cm,yunit=2.25cm,arrowsize=2pt 3}
\begin{pspicture*}(-3,-0.8)(2.4,2.6)
\psaxes[linewidth=1.25pt,Dx=10,Dy=10]{->}(0,0)(-3,-0.8)(2.4,2.6)
\psaxes[linewidth=1.25pt](0,0)(1,1)
\psplot[plotpoints=2000,linecolor=red,linewidth=1.25pt]{-3}{2.4}{2.71828 x neg exp 1 add ln}
\psplotTangent[linestyle=dashed]{0}{4}{2.71828 x neg exp 1 add ln}
\psline[linestyle=dashed](-3,2.7)(3.4,-0.45)
\psdots(-2.2,2.3)(2.3,0.1)
\uput[u](-2.2,2.3){$M_a$}\uput[u](2.3,0.1){$N_a$}
\uput[dl](-1.6,1.8){\red $\mathcal{C}$}\uput[d](2.3,-0.43){$T_0$}
\end{pspicture*}
\end{center}

\medskip

\begin{enumerate}
\item 
	\begin{enumerate}
		\item Déterminer la limite de la fonction $f$ en $-\infty$.
		\item Déterminer la limite de la fonction $f$ en $+ \infty$.
		
Interpréter graphiquement ce résultat.
		\item On admet que la fonction $f$ est dérivable sur $\R$ et on note $f'$ sa fonction dérivée.
		
Calculer $f'(x)$ puis montrer que, pour tout nombre réel $x$,\: $f'(x) = \dfrac{-1}{1 + \text{e}^x}$.
		\item Dresser le tableau de variations complet de la fonction $f$ sur $\R$.
	\end{enumerate}	
\item On note $T_0$ la tangente à la courbe $\mathcal{C}$ en son point d'abscisse $0$.
	\begin{enumerate}
		\item Déterminer une équation de la tangente $T_0$.
		\item Montrer que la fonction $f$ est convexe sur $\R$.
		\item En déduire que, pour tout nombre réel $x$, on a : 
		\[f(x) \geqslant  - \dfrac12x + \ln (2).\]
		\end{enumerate}	
\item Pour tout nombre réel $a$ différent de 0, on note $M_a$ et $N_a$ les points de la courbe $\mathcal{C}$ d'abscisses respectives $-a$ et $a$. 

On a donc : $M_a\left(-a~;~ f(-a)\right)$ et $N_a\left(a~;~ f(a)\right)$.
	\begin{enumerate}
		\item Montrer que, pour tout nombre réel $x$, on a : $f(x) - f(-x) = - x$.
		\item En déduire que les droites $T_0$ et $\left(M_aN_a\right)$ sont parallèles.
	\end{enumerate}
\end{enumerate}
